\subsection{终点误差控制的层次与条件}

全场最大含水率的阈值涉及时空最大值，单独给出相邻网格差不能认证连续方程的干燥终点。为区分计算误差、初值误差和模型误差，以下分别考察原全过程的数值对照、指定晚期初值的半离散误差管及连续标量屏障。三者的证明对象不同，不能互相替代。

\paragraph{全程细化与经验报告裕量。}

空间对照在同一容差下比较相邻网格；时间对照在同一最细网格下收紧容差。前3小时逐秒比较21个固定物理半径及真实表面，之后逐分钟比较；此外，对全程每分钟的完整径向节点并集比较分段线性重构。该并集上的最大差是两个声明重构在所保存分钟的精确空间最大差，仍不是连续PDE误差上界。三组参考计算的重放根、接受步哈希及原已存样本均逐位相同。

\begin{center}\small\begin{tabular}{lrr}\hline 项目 & 问题2--3 & 问题4 \\\hline

网格数 $N$ & 2560 & 5120 \\

同容差网格根差/s & 0.0536039 & 0.000469285 \\

同网格容差根差/s & 3.6048e-07 & 5.62928e-05 \\

条件半离散根区间宽/s & 2.76821e-06 & 0.00025197 \\

条件连续标量根区间宽/s & 0.102076 & 0.0035549 \\

\hline\end{tabular}\end{center}

表中后两行只针对下文重新指定的晚期初值问题。对全过程结果，报告裕量仍取完整同容差网格根差与同细网格容差根差之和，并向干侧取下一个0.0001小时刻度：

\[t_{\rm rep}=0.36\left\lceil\frac{t_{\rm fine}+|t_{N}-t_{N/2}|+|t_{N,\rm tight}-t_{N}|}{0.36}\right\rceil .\]

若恰好等于等号根，还需进至下一刻度。严格干燥由未舍入全节点最大含水率验证；四位小数显示可为0.1500。两问新报告时刻分别为 57.4796 h 和 51.0970 h。这一经验裕量不改称连续方程误差包围。独立8/16阶Gauss积分使用实际BDF稠密多项式的导数，核对水分边界通量及 $J\rho c_pT_t$ 与当前热边界通量，不将 $J\rho c_pT$ 的导数误作储热率。

\paragraph{半离散晚期初值与加权误差管。}

问题2--3在 206880 s、问题4在 183900 s重新指定旧网格检查点初值，其网格数分别为 1280 和 2560。温度精确等于恒定尾环境值。节点轨迹用独立Radau重启的三次多项式重构，末系数按端点差定义，保证时间接口在精确算术下连续。所有节点、参数和固定正权重的binary64数值作为精确实常数。

令上下候选为 $q\pm a(t)w$，其中保存的正节点权重 $w_i$ 不随时间改变。用向外区间给出全时间缺陷上界 $\eta\ge\max_i|q_i^{\prime}-f_i(q)|/w_i$，并沿同一个标量平移方向包围 $\mu\ge\max(0,\max_i\partial_\delta f_i(q+\delta w)/w_i)$。逐面核验管内水分同场非对角偏导非负，保留通量符号和外边界耗散，取

\[a(t_n+\tau)=(a_n+\eta\tau)e^{\mu\tau},\qquad 0\le\tau\le h_n.\]

温度盒通过恒环境下的离散热最大原理独立维持；不要求热湿联合系统合作。Cauchy余项控制扩散系数的时间级数尾项，44位Decimal向外算术覆盖全部时间区间。管闭合、阈值两侧包围与横截性分别核验；两问加权证书中的可能活动最大节点均为轴心，导数区间严格为负。该证书仍不包含零时刻至晚期初值的累积误差。

\paragraph{连续空间屏障及其明确初值。}

连续问题的水分初态另指定为通过上述旧检查点的唯一夹持 $C^2$ 径向三次样条，轴心导数为零，表面斜率精确满足非线性Robin条件。夹持样条方程严格对角占优，区间Thomas消元包住系数；精确解析表面斜率的级数余项经样条线性影响函数传播。相邻时间接口拥有同一节点值、同一半径和同一精确Robin斜率，故由样条唯一性得到整条空间函数及前两阶导数连续，不会遗漏内部界面通量源。

在 $x=r/R(t)$ 上，令 $L_t(C)=R^{-2}\{D(C)(C_{xx}+C_x/x)+D_C(C)C_x^2\}$，构造

\[U^{\pm}(t,x)=q(t,x)\pm a(t)e^{6(x^2-1)}.\]

这里权重是连续解析函数，与半离散证书的固定节点权重不同。全部空间单元和时间区间以直接剩余量及空间均值形式的交集包围，核验 $|q_t-L_t(q)|/w\le\eta$、沿屏障方向的导数上界及 $a<a_{\rm cap}$。首单元 $q=A_0+A_2v^2+A_3v^3$、$x=hv$，故 $\Delta q=(4A_2+9A_3v)/h^2$，轴心采用这一有限极限。表面还独立核验

\[\partial_\delta\{D(U)U_x+Rh_m(U-C_e)\}=D_C(U)(q_x+12\delta)+12D(U)+Rh_m\ge0.\]

两问实际晚期窗口的半径均正且恒定。正屏障把含水率限制在紧正区间；令 $u=F(C)$、$F^{\prime}(C)=D(C)>0$，可化为 $u_t=D(F^{-1}(u))R^{-2}\Delta u$，边界为 $\partial_nu+Rh_m(F^{-1}(u)-C_e)=0$。扩散一致为正，边界函数对 $u$ 的导数为 $Rh_m/D>0$；初态与Robin条件相容，径向轴心三次项仍属于 $C^{2,\alpha}$（$0<\alpha<1$），因此比较论证在该正则晚期问题中适用。上屏障径向不增、下屏障轴心值随时间不增均已区间核验，据此得到首次阈值穿越的两侧包围。

\begin{center}\small\begin{tabular}{ll}\hline 条件初值问题 & 等号根区间/s（显示端点向外取舍） \\\hline

问题2--3 半离散 & [206926.36143229, 206926.36143506] \\

问题2--3 连续标量 & [206926.31045170, 206926.41252820] \\

问题4 半离散 & [183948.83638749, 183948.83663947] \\

问题4 连续标量 & [183948.83473609, 183948.83829100] \\

\hline\end{tabular}\end{center}

\paragraph{新旧报告刻度查询与不可移用的结论。}

下表在各自连续晚期条件初值下查询原报告与新报告刻度。显示的上界均向外取舍，小于精确十进制阈值0.15；问题4两次报告刻度相同，仍保留两条来源。

\begin{center}\small\begin{tabular}{llrr}\hline 情景 & 刻度来源 & 物理时刻/s & 条件全域最大C上界 \\\hline

问题2--3 & 原报告 & 206926.92 & 0.149999852364934 \\

问题2--3 & 新报告 & 206926.56 & 0.149999957097033 \\

问题4 & 原报告 & 183949.2 & 0.149999814295733 \\

问题4 & 新报告 & 183949.2 & 0.149999814295733 \\

\hline\end{tabular}\end{center}

\textbf{新细网格全程等号根不必落入旧C2晚期初值区间。} 本次两问均如此：全过程从原始初态积分，连续证书则重新指定旧网格的晚期C2水分初态和精确恒温条件；二者继承误差和温度条件不同。上述条件刻度查询不能重命名为新细网格全过程的误差保证，也不能用区间未含新全程根来否定其对所指定条件问题的证明。

严格全程认证仍缺少原始连续初态至晚期检查点的误差管，以及非恒温时温度梯度误差和非合作耦合项的可验证传播界；仅有温差幅值界不足以控制 $D_T T_x C_x$。原热湿联合Jacobian存在负交叉偏导，不能直接套用合作系统比较原理。早期初边值不相容、分段环境、运动域、端面/运动学模型和输入不确定性也未由上述晚期证明覆盖。证书把冻结binary64参数及所计算尾均值视为精确实数，不包含十进制换算或尾均值计算舍入。因而现有结论是更充分的全程数值后验支持，加上两个严格限定初值的局部证明；尚非原全程连续模型的无条件终点认证。
